% =====================================================================
%  PARTIE 1 — De l'IoT à l'AIoT et au TinyML
% =====================================================================
\chapter{De l'IoT à l'AIoT et au TinyML}
\label{ch:fondements}

Avant d'étudier la manière dont un modèle d'intelligence artificielle peut être
exécuté sur un microcontrôleur, il est nécessaire de préciser le vocabulaire et
le cadre dans lequel s'inscrit ce travail. Cette partie définit successivement
l'objet connecté et l'architecture d'un système IoT
(section~\ref{sec:iot}), l'AIoT (section~\ref{sec:aiot}), le TinyML
(section~\ref{sec:tinyml}), puis compare les architectures centrées sur le cloud
à celles où l'inférence est réalisée dans l'objet
(section~\ref{sec:archi}).

% ---------------------------------------------------------------------
\section{L'Internet des objets}
\label{sec:iot}

\subsection{Qu'est-ce qu'un objet connecté ?}

L'expression \emph{Internet of Things} est attribuée à Kevin Ashton, qui
l'utilise dès 1999 pour décrire des objets capables de collecter eux-mêmes des
informations sur le monde physique, sans saisie humaine~\cite{ashton2009}. La
recommandation ITU-T~Y.2060 de l'Union internationale des télécommunications en
donne une définition de référence : l'IoT est \og une infrastructure mondiale
pour la société de l'information, qui permet de disposer de services évolués en
interconnectant des objets (physiques ou virtuels) grâce aux technologies de
l'information et de la communication interopérables existantes ou en
évolution\fg{}~\cite{itu2012y2060}.

Un \textbf{objet connecté} est donc un dispositif physique qui associe :
\begin{itemize}
  \item un ou plusieurs \textbf{capteurs} qui mesurent une grandeur physique
        (température, accélération, son, lumière, humidité, pression\dots) ;
  \item une \textbf{unité de traitement}, le plus souvent un microcontrôleur
        (MCU), qui pilote les capteurs et met en forme les données ;
  \item un \textbf{module de communication} (Wi-Fi, Bluetooth Low Energy,
        LoRaWAN, NB-IoT, Zigbee\dots) qui relie l'objet au réseau ;
  \item éventuellement des \textbf{actionneurs} (relais, moteur, LED,
        buzzer) qui agissent sur l'environnement ;
  \item une \textbf{source d'énergie}, souvent une batterie ou une pile, parfois
        complétée par de la récupération d'énergie (panneau solaire, vibrations).
\end{itemize}

La carte ESP32 d'Espressif illustre bien ce type de plateforme : un
microcontrôleur cadencé à 240~MHz, environ 520~Ko de SRAM, du Wi-Fi et du
Bluetooth intégrés, pour un coût de quelques euros~\cite{espressif-esp32}.
La contrainte majeure de ces objets n'est pas tant leur puissance de calcul que
leur \textbf{budget énergétique} : un capteur installé dans un champ ou sur une
machine doit souvent fonctionner plusieurs mois, voire plusieurs années, sans
intervention.

\subsection{Architecture générale d'un système IoT}

Il n'existe pas d'architecture IoT unique, mais la littérature converge vers un
découpage en couches~\cite{alfuqaha2015,atzori2010}. Le modèle le plus courant
distingue quatre niveaux, représentés à la figure~\ref{fig:archi-iot} :

\begin{enumerate}
  \item \textbf{La couche perception} (ou couche objets) regroupe les capteurs et
        actionneurs. Elle convertit des phénomènes physiques en signaux
        numériques et, inversement, exécute des commandes.
  \item \textbf{La couche réseau} (ou transport) achemine les données. Elle
        s'appuie sur des protocoles radio adaptés aux faibles débits et à la
        faible consommation (BLE, Zigbee, LoRaWAN, NB-IoT) ainsi que sur des
        passerelles (\emph{gateways}) qui relaient les messages vers Internet.
  \item \textbf{La couche traitement} (ou plateforme, \emph{middleware}) stocke,
        agrège et analyse les données, le plus souvent dans le cloud. On y trouve
        les bases de données, les tableaux de bord et, aujourd'hui, les modèles
        d'apprentissage automatique.
  \item \textbf{La couche application} fournit le service final à l'utilisateur :
        application mobile, alerte de maintenance, pilotage d'un système
        d'irrigation, etc.
\end{enumerate}

\begin{figure}[H]
\centering
\begin{tikzpicture}[
    couche/.style={draw=#1, fill=#1!12, rounded corners=3pt, minimum width=13.6cm,
                   minimum height=1.25cm},
    lbl/.style={font=\bfseries\small, anchor=west, text width=3cm},
    ex/.style={font=\small, anchor=west, text width=9.6cm}]
  \node[couche=orange] (app) at (0,4.5) {};
  \node[lbl] at ($(app.west)+(0.2,0)$) {4. Application};
  \node[ex]  at ($(app.west)+(3.6,0)$) {Tableaux de bord, applications mobiles, alertes,
       maintenance, villes et maisons intelligentes};
  \node[couche=vert] (trt) at (0,3) {};
  \node[lbl] at ($(trt.west)+(0.2,0)$) {3. Traitement};
  \node[ex]  at ($(trt.west)+(3.6,0)$) {Cloud / serveurs : stockage, base de données,
       analyse, entraînement des modèles d'IA};
  \node[couche=bleu] (res) at (0,1.5) {};
  \node[lbl] at ($(res.west)+(0.2,0)$) {2. Réseau};
  \node[ex]  at ($(res.west)+(3.6,0)$) {Wi-Fi, BLE, Zigbee, LoRaWAN, NB-IoT, passerelles ;
       protocoles MQTT, CoAP, HTTP};
  \node[couche=gray] (per) at (0,0) {};
  \node[lbl] at ($(per.west)+(0.2,0)$) {1. Perception};
  \node[ex]  at ($(per.west)+(3.6,0)$) {Capteurs (température, accéléromètre, micro\dots),
       actionneurs, microcontrôleurs};
  \draw[fleche] ($(per.east)+(0.3,0)$) -- node[right,font=\footnotesize,align=left]{données\\montantes} ($(app.east)+(0.3,0)$);
\end{tikzpicture}
\caption{Architecture en quatre couches d'un système IoT (d'après~\cite{alfuqaha2015}).}
\label{fig:archi-iot}
\end{figure}

Dans ce schéma classique, l'« intelligence » du système réside dans la couche
traitement. L'objet est un simple producteur de données : il mesure, transmet et
attend éventuellement un ordre en retour. Cette organisation a l'avantage de la
simplicité et de la puissance (le cloud dispose de ressources pratiquement
illimitées), mais elle fait dépendre toute décision de la disponibilité et de la
qualité du réseau.

\subsection{Les limites du modèle « tout cloud »}

Plusieurs difficultés apparaissent lorsque le nombre d'objets et le volume de
données augmentent~\cite{shi2016edge} :
\begin{itemize}
  \item \textbf{La bande passante} : un accéléromètre échantillonné à 1~kHz sur
        trois axes en 16~bits produit environ 6~Ko/s, soit plus de 500~Mo par
        jour et par capteur. Multiplier ce flux par des centaines de machines
        sature rapidement un réseau LoRaWAN ou NB-IoT, dont les débits se
        comptent en kilobits par seconde.
  \item \textbf{La latence} : l'aller-retour vers un serveur ajoute des délais
        variables, parfois de plusieurs centaines de millisecondes, incompatibles
        avec certaines réactions (arrêt d'une machine, détection d'une chute).
  \item \textbf{L'énergie} : sur un objet sur batterie, la transmission radio
        est souvent le poste de consommation le plus important, bien devant le
        calcul~\cite{warden2019tinyml}.
  \item \textbf{La confidentialité} : envoyer en continu le son d'une pièce ou
        le rythme cardiaque d'une personne multiplie les risques de fuite.
  \item \textbf{La disponibilité} : sans connexion, l'objet devient « aveugle ».
\end{itemize}
Ces limites motivent le déplacement d'une partie de l'analyse vers la
périphérie du réseau (\emph{edge}), puis vers l'objet lui-même.

% ---------------------------------------------------------------------
\section{L'AIoT : l'intelligence artificielle au service des objets}
\label{sec:aiot}

\subsection{Définition}

L'\textbf{AIoT} (\emph{Artificial Intelligence of Things}) désigne la
combinaison des technologies de l'intelligence artificielle, en particulier de
l'apprentissage automatique, avec l'infrastructure de l'Internet des
objets~\cite{zhang2021aiot}. L'objectif est de transformer les données
collectées par les objets en informations exploitables (reconnaître une
situation, détecter une anomalie, anticiper un événement) et, idéalement, de
déclencher une action sans intervention humaine.

L'IA peut être placée à différents endroits de l'architecture : dans le
cloud, sur une passerelle ou un serveur de proximité (\emph{edge server}), ou
directement dans l'objet (\emph{on-device AI}). L'AIoT recouvre l'ensemble de
ces configurations ; ce qui la caractérise est le passage d'objets qui
\emph{mesurent} à des systèmes qui \emph{comprennent} et \emph{décident}.

\subsection{Pourquoi ajouter de l'IA aux objets connectés ?}

Les règles fixes (\og si la température dépasse 80~°C, alors alarme\fg{})
suffisent pour des phénomènes simples. Elles deviennent inadaptées lorsque le
signal est complexe : reconnaître un mot prononcé, distinguer le bruit normal
d'un roulement de celui d'un roulement usé, identifier un geste à partir d'un
accéléromètre. Dans ces cas, il est très difficile d'écrire les règles à la
main, alors qu'un modèle d'apprentissage automatique peut les apprendre à
partir d'exemples. L'IA apporte ainsi aux systèmes IoT :
\begin{itemize}
  \item la \textbf{reconnaissance de motifs} dans des signaux bruités (audio,
        vibrations, images) ;
  \item la \textbf{détection d'anomalies}, c'est-à-dire l'identification
        d'écarts par rapport à un comportement appris comme normal ;
  \item la \textbf{prédiction} de l'évolution d'une grandeur (consommation,
        usure, rendement) ;
  \item l'\textbf{adaptation} au contexte et à l'utilisateur.
\end{itemize}

\subsection{Exemples d'utilisation}

\begin{table}[H]
\centering
\small
\begin{tabularx}{\linewidth}{>{\bfseries}l X X}
\toprule
Domaine & Données collectées & Apport de l'IA \\
\midrule
Industrie      & Vibrations, courant, température des moteurs & Détection d'anomalies, maintenance prédictive \\
Maison         & Son, présence, consommation électrique & Reconnaissance de mots-clés, automatisation \\
Agriculture    & Humidité du sol, météo, images de feuilles & Irrigation optimisée, détection de maladies \\
Santé          & Fréquence cardiaque, mouvements & Détection de chutes ou d'arythmies \\
Transport      & Accélérations, images, données moteur & Aide à la conduite, surveillance du véhicule \\
Environnement  & Sons de la forêt, qualité de l'air & Détection de tronçonneuses, alerte pollution \\
\bottomrule
\end{tabularx}
\caption{Exemples d'applications de l'AIoT.}
\label{tab:exemples-aiot}
\end{table}

Ces exemples, développés dans la partie~\ref{ch:benefices}, ont un point
commun : la décision utile (\og anomalie\fg{}, \og mot-clé reconnu\fg{},
\og chute\fg{}) représente une quantité d'information infime par rapport aux
données brutes dont elle est issue. C'est précisément ce qui rend intéressante
l'exécution du modèle au plus près du capteur.

% ---------------------------------------------------------------------
\newpage
\section{Le TinyML : l'apprentissage automatique sur microcontrôleur}
\label{sec:tinyml}

\subsection{Définition}

Le \textbf{TinyML} (\emph{Tiny Machine Learning}) désigne le domaine qui
regroupe les technologies, matérielles et logicielles, permettant d'exécuter
des modèles d'apprentissage automatique sur des appareils à très faibles
ressources, typiquement des microcontrôleurs, avec une consommation de l'ordre
du milliwatt ou moins~\cite{warden2019tinyml,banbury2020benchmarking}. Le
terme a été popularisé en 2019 par l'ouvrage de Pete Warden et Daniel
Situnayake~\cite{warden2019tinyml} et par la création de la communauté TinyML
Foundation, devenue depuis EDGE~AI FOUNDATION~\cite{tinymlfoundation}.

Le TinyML se situe à l'extrémité du continuum de l'IA embarquée, comme le
montre le tableau~\ref{tab:spectre}. Là où un serveur dispose de dizaines de
gigaoctets de mémoire et un smartphone de plusieurs gigaoctets, un
microcontrôleur n'en possède que quelques centaines de kilooctets, soit un
facteur de l'ordre de $10^5$.

\begin{table}[H]
\centering
\small
\begin{tabularx}{\linewidth}{l X X X}
\toprule
 & \textbf{Cloud AI} & \textbf{Mobile / Edge AI} & \textbf{TinyML} \\
\midrule
Matériel        & Serveurs GPU/TPU & Smartphone, Raspberry Pi, Jetson & Microcontrôleur (ESP32, Arm Cortex-M) \\
Mémoire vive    & Dizaines à centaines de Go & 1 à 16 Go & 32 Ko à quelques Mo \\
Stockage        & To & Dizaines de Go & 256 Ko à quelques Mo de flash \\
Puissance       & Centaines de watts & Quelques watts & Milliwatts \\
Système         & Linux, conteneurs & Android, Linux & Pas d'OS ou RTOS (FreeRTOS) \\
Coût unitaire   & Milliers d'euros & 50 à 500 € & 1 à 10 € \\
Taille de modèle & Go & Dizaines de Mo & Ko à quelques centaines de Ko \\
\bottomrule
\end{tabularx}
\caption{Positionnement du TinyML par rapport aux autres formes d'IA (ordres de grandeur, d'après~\cite{banbury2020benchmarking,lin2020mcunet}).}
\label{tab:spectre}
\end{table}

\subsection{Les contraintes du TinyML}

Concevoir une application TinyML revient à respecter simultanément plusieurs
contraintes très fortes~\cite{banbury2020benchmarking,ray2022review} :

\paragraph{La mémoire.} Un microcontrôleur dispose de deux types de mémoire.
La \textbf{mémoire flash} (non volatile) stocke le programme et les poids du
modèle ; la \textbf{SRAM} (volatile) contient les données de travail, en
particulier les activations intermédiaires calculées pendant l'inférence. Les
deux sont limitées : un modèle doit typiquement tenir dans quelques centaines de
kilooctets de flash et ses activations dans quelques dizaines à centaines de
kilooctets de SRAM. Lin \emph{et al.} montrent que c'est souvent la SRAM, plus
encore que la flash, qui constitue le goulot d'étranglement~\cite{lin2020mcunet}.

\paragraph{La puissance de calcul.} Les microcontrôleurs fonctionnent à des
fréquences de quelques dizaines à quelques centaines de mégahertz, souvent sans
unité de calcul flottant performante. Un réseau de neurones de quelques
millions d'opérations est envisageable ; un modèle de plusieurs milliards
d'opérations ne l'est pas.

\paragraph{L'énergie.} Pour fonctionner des mois sur une pile, l'objet doit
passer l'essentiel de son temps en veille et limiter la durée de chaque
inférence. La consommation se mesure alors en milliwatts et l'énergie par
inférence en microjoules ou millijoules~\cite{banbury2021mlperftiny}.

\paragraph{L'environnement logiciel.} Pas de système d'exploitation complet,
pas toujours d'allocation dynamique de mémoire, pas de Python : le modèle est
compilé avec le micrologiciel (\emph{firmware}) en C/C++ et exécuté par un
moteur d'inférence minimaliste comme TensorFlow Lite for Microcontrollers
(TFLM)~\cite{david2021tflm}.

\subsection{Entraînement et inférence : une séparation fondamentale}

Il est essentiel de distinguer les deux phases de vie d'un modèle :
\begin{itemize}
  \item l'\textbf{entraînement}, qui consiste à ajuster les paramètres du modèle
        à partir d'un grand nombre d'exemples. Il demande beaucoup de mémoire
        et de calcul et s'effectue sur un ordinateur ou dans le cloud ;
  \item l'\textbf{inférence}, qui consiste à appliquer le modèle déjà entraîné à
        une nouvelle donnée pour obtenir une prédiction. C'est la seule phase
        exécutée sur le microcontrôleur dans la grande majorité des systèmes
        TinyML~\cite{david2021tflm}.
\end{itemize}
Entre les deux, le modèle doit être \textbf{optimisé} (réduit, quantifié) puis
\textbf{converti} dans un format exécutable par le microcontrôleur. Ces étapes
sont détaillées dans la partie~\ref{ch:principes}.

\begin{aretenir}
L'IoT connecte les objets, l'AIoT leur donne la capacité d'interpréter leurs
données, et le TinyML pousse cette logique jusqu'aux plus petits composants :
des microcontrôleurs à quelques euros, alimentés sur pile, qui exécutent
localement un modèle entraîné ailleurs.
\end{aretenir}

% ---------------------------------------------------------------------
\section{Fonctionnement d'une architecture AIoT/TinyML}
\label{sec:archi}

\subsection{La chaîne capteur -- microcontrôleur -- modèle -- action}

Dans un système TinyML, le traitement suit une boucle locale composée de cinq
étapes (figure~\ref{fig:boucle}) :

\begin{enumerate}
  \item \textbf{Acquisition} : le capteur mesure le phénomène (par exemple
        l'accélération d'un moteur) et le microcontrôleur lit les échantillons
        via un bus I\textsuperscript{2}C, SPI ou un convertisseur
        analogique-numérique.
  \item \textbf{Prétraitement} : les échantillons sont regroupés dans une
        fenêtre temporelle (par exemple une seconde) puis transformés en
        caractéristiques plus compactes : moyenne, écart-type, valeur efficace,
        spectre par transformée de Fourier (FFT), coefficients MFCC pour l'audio.
  \item \textbf{Inférence} : le vecteur de caractéristiques est fourni au modèle,
        qui renvoie une prédiction (classe, score d'anomalie, valeur).
  \item \textbf{Décision} : la prédiction est comparée à un seuil ou interprétée
        par une règle simple.
  \item \textbf{Action et communication} : l'objet agit (LED, relais, arrêt
        d'urgence) et ne transmet au réseau que l'information utile, par exemple
        un message MQTT \og anomalie détectée\fg{}.
\end{enumerate}

\begin{figure}[H]
\centering
\begin{tikzpicture}[node distance=0.55cm]
  \node[bloc, minimum width=2.2cm] (cap) {Capteur\\\footnotesize(accéléromètre)};
  \node[bloc, minimum width=2.2cm, right=of cap] (pre) {Prétraitement\\\footnotesize(fenêtre, FFT)};
  \node[blocv, minimum width=2.2cm, right=of pre] (inf) {Inférence\\\footnotesize(modèle TinyML)};
  \node[blocv, minimum width=2.0cm, right=of inf] (dec) {Décision\\\footnotesize(seuil)};
  \node[bloco, minimum width=2.0cm, right=of dec] (act) {Action\\\footnotesize(alerte, LED)};
  \draw[fleche] (cap) -- (pre);
  \draw[fleche] (pre) -- (inf);
  \draw[fleche] (inf) -- (dec);
  \draw[fleche] (dec) -- (act);
  \begin{scope}[on background layer]
    \node[draw=gris, dashed, rounded corners=6pt, fit=(pre)(inf)(dec), inner sep=8pt,
          label={[font=\footnotesize\itshape,text=gris]below:dans le microcontrôleur (ESP32)}] {};
  \end{scope}
  \node[bloco, minimum width=2.6cm, above=1.1cm of act, xshift=-1.2cm] (cloud) {Cloud / tableau de bord\\\footnotesize(alertes uniquement)};
  \draw[fleche, dashed] (act) -- node[right,font=\footnotesize]{MQTT} (cloud);
\end{tikzpicture}
\caption{Boucle de traitement locale d'un objet AIoT/TinyML.}
\label{fig:boucle}
\end{figure}

Le listing~\ref{lst:boucle} montre, sous forme simplifiée, la structure du
programme principal d'un ESP32 qui applique ce principe. Le modèle est
interrogé à chaque fenêtre de mesures et seule une alerte est publiée.

\begin{lstlisting}[language=C++, caption={Structure simplifiée de la boucle d'inférence sur ESP32.}, label={lst:boucle}]
void loop() {
  // 1. Acquisition : une fenêtre de 1 s à 1 kHz sur 3 axes
  lire_accelerometre(fenetre, 1000);

  // 2. Prétraitement : extraction de caractéristiques (RMS, FFT...)
  calculer_caracteristiques(fenetre, entree_modele);

  // 3. Inférence locale (TensorFlow Lite for Microcontrollers)
  interpreter->Invoke();
  float score = sortie_modele[0];

  // 4-5. Décision et action : on ne transmet que l'événement utile
  if (score > SEUIL_ANOMALIE) {
    digitalWrite(LED_ALERTE, HIGH);
    mqtt.publish("usine/moteur1/alerte", "anomalie");
  }
  esp_light_sleep_start();   // mise en veille jusqu'à la prochaine fenêtre
}
\end{lstlisting}

\subsection{Comparaison avec une architecture centrée sur le cloud}

Dans une architecture IoT classique, l'objet transmet en continu ses mesures
brutes ; la passerelle les relaie vers le cloud, où le modèle est exécuté avant
que la décision ne redescende éventuellement vers l'objet. Dans une architecture
TinyML, l'inférence est réalisée dans l'objet et seules les décisions remontent.
La figure~\ref{fig:cloud-vs-edge} compare les deux flux.

\begin{figure}[H]
\centering
\begin{tikzpicture}[node distance=0.9cm]
  % ---- Architecture cloud
  \node[font=\bfseries\small] (t1) at (0,0) {a) Architecture centrée sur le cloud};
  \node[bloc, below=0.4cm of t1, xshift=-5cm, minimum width=2cm] (c1) {Capteur\\+ MCU};
  \node[bloc, right=of c1, minimum width=2cm] (g1) {Passerelle};
  \node[bloco, right=of g1, minimum width=2.4cm] (cl1) {Cloud\\(modèle IA)};
  \node[bloc, right=of cl1, minimum width=2cm] (a1) {Décision\\/ action};
  \draw[fleche, draw=orange, line width=1.6pt] (c1) -- node[above,font=\scriptsize]{brut} (g1);
  \draw[fleche, draw=orange, line width=1.6pt] (g1) -- node[above,font=\scriptsize]{brut} (cl1);
  \draw[fleche] (cl1) -- (a1);
  \draw[fleche, dashed] (a1.south) -- ++(0,-0.45) -| node[pos=0.25,below,font=\scriptsize]{commande en retour (latence réseau)} (c1.south);

  % ---- Architecture TinyML
  \node[font=\bfseries\small, below=2.4cm of t1] (t2) {b) Architecture TinyML (inférence locale)};
  \node[blocv, below=0.4cm of t2, xshift=-5cm, minimum width=2cm] (c2) {Capteur\\+ MCU\\+ modèle};
  \node[blocv, right=of c2, minimum width=2cm] (a2) {Décision\\/ action\\locale};
  \node[bloc, right=of a2, minimum width=2cm] (g2) {Passerelle};
  \node[bloc, right=of g2, minimum width=2.4cm] (cl2) {Cloud\\(supervision)};
  \draw[fleche] (c2) -- (a2);
  \draw[fleche, draw=vert, line width=0.6pt] (a2) -- node[above,font=\scriptsize]{alerte} (g2);
  \draw[fleche, draw=vert, line width=0.6pt] (g2) -- node[above,font=\scriptsize]{alerte} (cl2);
\end{tikzpicture}
\caption{Flux de données : architecture cloud (données brutes, flèches épaisses) et architecture TinyML (seules les alertes remontent).}
\label{fig:cloud-vs-edge}
\end{figure}

\begin{table}[H]
\centering
\small
\begin{tabularx}{\linewidth}{>{\bfseries}l X X}
\toprule
Critère & Architecture cloud & Architecture TinyML \\
\midrule
Lieu de l'inférence      & Serveur distant & Microcontrôleur de l'objet \\
Données transmises       & Mesures brutes, en continu & Résultats et alertes uniquement \\
Latence de décision      & Dépend du réseau (de dizaines de ms à plusieurs secondes) & Quelques ms à quelques dizaines de ms, indépendante du réseau \\
Fonctionnement hors ligne & Impossible & Possible pour les fonctions locales \\
Confidentialité          & Données sensibles transférées et stockées à distance & Données brutes conservées dans l'objet \\
Complexité des modèles   & Pratiquement illimitée & Fortement contrainte (Ko de mémoire) \\
Mise à jour du modèle    & Simple (côté serveur) & Plus délicate (mise à jour du firmware, OTA) \\
Coût récurrent           & Abonnement cloud, transfert de données & Faible, surtout lié à la supervision \\
\bottomrule
\end{tabularx}
\caption{Comparaison entre une architecture centrée sur le cloud et une architecture TinyML.}
\label{tab:cloud-vs-tinyml}
\end{table}

Les deux approches ne s'opposent pas nécessairement. Dans la pratique, les
architectures AIoT sont souvent \textbf{hybrides} : l'objet filtre et
interprète localement les données, une passerelle ou le cloud agrège les
résultats de nombreux objets, réalise les analyses lourdes et ré-entraîne
périodiquement les modèles, qui sont ensuite redéployés dans les objets. Cette
complémentarité est discutée dans la partie~\ref{ch:limites}.

\subsection{Le cas d'étude retenu}

Pour illustrer concrètement ces principes, ce rapport s'appuie sur un exemple
fil rouge : la \textbf{détection d'anomalies vibratoires sur un moteur}. Un
accéléromètre (de type MPU6050 ou ADXL345) est fixé sur le carter du moteur et
relié à un ESP32. Un petit modèle, entraîné sur ordinateur à partir de mesures
en fonctionnement normal, est quantifié puis intégré au firmware. L'ESP32 analyse
chaque seconde de vibrations et publie une alerte lorsque le comportement
s'écarte de la normale. Ce scénario mobilise toutes les briques étudiées dans
les parties suivantes : pipeline de développement, compression du modèle,
matériel et outils, bénéfices et limites.
